% !TeX root = ../../金融经济学_王江_讲义及习题.tex
% 03_Arrow-Debreu经济 课后习题（由章节文件拆分）

    \begin{problem}
        \problemno{例 3.3} 考虑“Lucas 树经济”，1 期有两个概率相等的状态 $a$ 和 $b$。假设 1 期的两个可能状态的状态价格分别为 $\phi_{a}$ 和 $\phi_{b}$。考虑一个参与者，他的禀赋为 $(e_{0};e_{1a};e_{1b})$。他具有如下对数形式的效用函数：
        \[
            U(c_{0},c_{1a},c_{1b})=\log c_{0}+\tfrac12(\log c_{1a}+\log c_{1b})
        \]
        
        他的最优消费/组合选择是什么？
        \end{problem}
        \begin{solution}
        给定状态价格和他的禀赋，他的总财富是 $w=e_{0}+\phi_{a}e_{1a}+\phi_{b}e_{1b}$。他的最优化问题是
        \[
            \begin{aligned}
                \max_{\{c_{0},c_{1a},c_{1b}\}}\ &\log c_{0}+\tfrac12(\log c_{1a}+\log c_{1b})\\
                \mathrm{s.t.}\quad &w-(c_{0}+\phi_{a}c_{1a}+\phi_{b}c_{1b})=0\\
                &c_{0},c_{1a},c_{1b}>0
            \end{aligned}
        \]

        假设 $w>0$、$\phi_a,\phi_b>0$。对数效用的消费域为 $c_0,c_{1a},c_{1b}>0$；最优解在内部，可直接使用等式预算的拉格朗日函数
        \[
            \mathcal{L}=\log c_{0}+\tfrac12(\log c_{1a}+\log c_{1b})+\lambda\bigl(w-c_0-\phi_a c_{1a}-\phi_b c_{1b}\bigr).
        \]
        一阶条件为
        \[
          \frac1{c_0}=\lambda,\qquad
          \frac1{2c_{1a}}=\lambda\phi_a,\qquad
          \frac1{2c_{1b}}=\lambda\phi_b,\qquad
          c_0+\phi_a c_{1a}+\phi_b c_{1b}=w.
        \]
        当消费趋近 $0$ 时，对数效用趋于负无穷、边际效用趋于正无穷，因此最优消费在每一时期、每一状态都严格为正。由上述一阶条件得到最优消费：
        \[
            c_{0}^{*}=\frac{1}{\lambda},\qquad
            c_{1a}^{*}=\frac{1}{\lambda}\frac{1}{2\phi_{a}},\qquad
            c_{1b}^{*}=\frac{1}{\lambda}\frac{1}{2\phi_{b}}
        \]

        把 $c^{*}$ 代入预算约束，我们可以得到 $\lambda$ 的解：
        \[
            \frac{1}{\lambda}+\frac{1}{2\lambda}+\frac{1}{2\lambda}=w\qquad\Rightarrow\qquad\lambda=\frac{2}{w}
        \]

        最后，我们有
        \[
            c_{0}^{*}=\frac12 w,\qquad
            c_{1a}^{*}=\frac14\frac{w}{\phi_{a}},\qquad
            c_{1b}^{*}=\frac14\frac{w}{\phi_{b}}
        \]

        可以看出，参与者把一半财富用作现在的消费，把另外一半财富作为未来的消费。在财富给定时，某一状态下的消费与对应状态价格成反比；若价格变化也改变禀赋市值，还须计入财富效应。这个结果很直观。状态价格高的状态下的消费更昂贵。结果，参与者在这些状态下选择较低的消费。
        \end{solution}

    \begin{problem}
        \problemno{例 3.4} 考虑一个经济，在 1 期有两个概率相等的状态 $a$ 和 $b$。经济中有参与者 1 和 2，他们具有的禀赋分别为：
        \begin{center}
        \WJEventTree[left={$e_{1}{:}100$}]{0,0}
        \qquad
        \WJEventTree[left={$e_{2}{:}0$}]{200,50}
        \end{center}
        两个参与者都具有如下形式的对数效用函数：
        \[
            U(c)=\log c_{0}+\tfrac12(\log c_{a}+\log c_{b})
        \]
        \end{problem}
        \begin{solution}
        在市场上存在一组完全的状态或有证券可以交易。因为有两个状态，因而只有两个状态或有证券。
        
        现在我们开始分析这个经济的均衡。从给定交易证券价格下参与者的最优化问题开始。记 $\phi=\begin{bmatrix}\phi_{a} \\ \phi_{b}\end{bmatrix}$ 为状态价格（向量），即两个状态或有证券的价格。我们可以定义每个参与者的财富为 $w=\hat{\phi}^{\top}e$，这里 $\hat{\phi}=\begin{bmatrix}1 \\ \phi\end{bmatrix}$，而 $e$ 是他的禀赋。这时，最优化问题变成了：
        \[
            \begin{aligned}
                \max_{c}\ &\log c_{0}+\tfrac12(\log c_{a}+\log c_{b})\\
                \mathrm{s.t.}\quad &c_{0}+\phi_{a}c_{a}+\phi_{b}c_{b}=w
            \end{aligned}
        \]

        由例 3.3，最优消费为
        \[
            c_{k,0}^{*}=\tfrac12 w_{k},\qquad
            c_{k,a}^{*}=\tfrac14\frac{w_{k}}{\phi_{a}},\qquad
            c_{k,b}^{*}=\tfrac14\frac{w_{k}}{\phi_{b}},\qquad k=1,2
        \]
        在状态价格 $\phi_{a},\phi_{b}$ 下，参与者按禀赋市值计的财富为 $w_{k}=e_{k,0}+\phi_{a}e_{k,a}+\phi_{b}e_{k,b}$（0 期消费品价格为 $1$）。代入禀赋 $(100;0,0)$ 与 $(0;200,50)$：
        \[
            w_{1}=100+0\cdot\phi_{a}+0\cdot\phi_{b}=100,\qquad
            w_{2}=0+200\phi_{a}+50\phi_{b}=200\phi_{a}+50\phi_{b}.
        \]

        \textbf{0 期商品市场出清}（两人 0 期消费之和等于经济 0 期总禀赋 $100$）。由例 3.3，$c_{k,0}^{*}=\tfrac12 w_{k}$，故
        \[
            c_{1,0}^{*}+c_{2,0}^{*}=\tfrac12 w_{1}+\tfrac12 w_{2}=\tfrac12(w_{1}+w_{2})=100
            \quad\Longrightarrow\quad w_{1}+w_{2}=200.
        \]

        把 $w_{1},w_{2}$ 代入 $w_{1}+w_{2}=200$：
        \[
            100+(200\phi_{a}+50\phi_{b})=200.
        \]

        移项得 $200\phi_{a}+50\phi_{b}=100$；两边同除以 $50$，得到与状态价格有关的第一个方程
        \[
            4\phi_{a}+\phi_{b}=2.
        \]

        \textbf{1 期·状态 $a$ 商品市场出清}（状态 $a$ 下总消费 $200$）。由 $c_{k,a}^{*}=\tfrac14 w_{k}/\phi_{a}$，
        \[
            c_{1,a}^{*}+c_{2,a}^{*}=\tfrac{w_{1}}{4\phi_{a}}+\tfrac{w_{2}}{4\phi_{a}}
            =\frac{w_{1}+w_{2}}{4\phi_{a}}=200
            \quad\Longrightarrow\quad w_{1}+w_{2}=800\phi_{a}.
        \]

        与 $w_{1}+w_{2}=200$ 联立得 $\phi_{a}=\dfrac14$，从而 $\phi_{b}=1$。状态 $b$ 出清 $\dfrac{w_{1}+w_{2}}{4\phi_{b}}=50$ 自动成立。        此时 $w_{1}=w_{2}=100$，最优消费为
        \[
            \begin{aligned}
                c_{1,0}^{*}=c_{2,0}^{*}&=50\\
                c_{1,a}^{*}=c_{2,a}^{*}&=100\\
                c_{1,b}^{*}=c_{2,b}^{*}&=25
            \end{aligned}
        \]
        \end{solution}

\begin{problem}
\problemno{3.1*} 考虑如下经济，在 1 期有两个可能状态 $a$ 和 $b$：
\begin{center}
\WJEventTree{a,b}
\end{center}
\begin{enumerate}[label=(\alph*)]
    \item 描述所有 Arrow-Debreu 证券的支付向量。记这些证券的价格向量为 $\phi$。
    \item 考虑一个拥有如下禀赋的参与者：
    \begin{center}
    \WJEventTree[left={$0$}]{2,1}
    \end{center}
    把他的禀赋表示成 Arrow-Debreu 证券的组合。
    \item 计算他的金融财富。写出他的预算集。
    \item 假设参与者的效用函数如下：
    \[
        U(c_{0},c_{1a},c_{1b})=-e^{-c_{0}}-\tfrac12\bigl(e^{-c_{1a}}+e^{-c_{1b}}\bigr).
    \]
    不考虑消费的非负约束，写出他的优化问题。求解他的最优消费选择。
    \item 讨论他的消费如何依赖于 Arrow-Debreu 证券的价格向量 $\phi$。
    \item 证明在某些价格下，他（在某些时期/状态下）的消费可能是负的。
\end{enumerate}
\end{problem}
\begin{solution}
\begin{enumerate}[label=(\alph*)]
    \item Arrow-Debreu 证券的支付向量是 $X_{a}=\begin{bmatrix}1 \\ 0\end{bmatrix}$，$X_{b}=\begin{bmatrix}0 \\ 1\end{bmatrix}$；
    \item $2X_{a}+X_{b}$；
    \item 0 期禀赋为 $0$，1 期禀赋是 $2$ 单位状态 $a$ 或有要求与 $1$ 单位状态 $b$ 或有要求，金融财富为 $w=2\phi_{a}+\phi_{b}$。在消费集 $\mathrm{R}_{+}^{3}$ 上，预算集是
    \[
        \bigl\{c\in\mathrm{R}_{+}^{3}:c_{0}+\phi_{a}c_{1a}+\phi_{b}c_{1b}=w\bigr\}.
    \]
    \item 本问不施加 $c\ge 0$，在 $\phi_{a}>0$、$\phi_{b}>0$ 时于 $c\in\mathrm{R}^{3}$ 上求解
    \[
        \begin{aligned}
            \max_{\{c_{0},c_{1a},c_{1b}\}}\ &\bigl\{-e^{-c_{0}}-\tfrac12\bigl(e^{-c_{1a}}+e^{-c_{1b}}\bigr)\bigr\}\\
            \mathrm{s.t.}\quad &c_{0}+\phi_{a}c_{1a}+\phi_{b}c_{1b}=w=2\phi_{a}+\phi_{b}
        \end{aligned}
    \]
    拉格朗日函数为
    \[
        \mathcal{L}=-e^{-c_{0}}-\tfrac12\bigl(e^{-c_{1a}}+e^{-c_{1b}}\bigr)+\lambda\bigl(w-c_{0}-\phi_{a}c_{1a}-\phi_{b}c_{1b}\bigr)
    \]
    对消费求导，一阶条件为
    \[
        e^{-c_{0}}=\lambda,\qquad
        \tfrac12 e^{-c_{1a}}=\lambda\phi_{a},\qquad
        \tfrac12 e^{-c_{1b}}=\lambda\phi_{b}
    \]
    从而最优消费为
    \[
        c_{0}^{*}=-\log\lambda,\qquad
        c_{1a}^{*}=-\log(2\lambda\phi_{a}),\qquad
        c_{1b}^{*}=-\log(2\lambda\phi_{b})
    \]
    代入预算约束：
    \[
        -\log\lambda-\phi_{a}\log(2\lambda\phi_{a})-\phi_{b}\log(2\lambda\phi_{b})=w
    \]
    即 $-(1+\phi_{a}+\phi_{b})\log\lambda-\phi_{a}\log(2\phi_{a})-\phi_{b}\log(2\phi_{b})=w$。记 $S=1+\phi_{a}+\phi_{b}$，则
    \[
        \log\lambda=-\frac{w+\phi_{a}\log(2\phi_{a})+\phi_{b}\log(2\phi_{b})}{S}
    \]
    从而 $c_{0}^{*}=-\log\lambda$，且 $c_{1a}^{*}=-\log\lambda-\log(2\phi_{a})$、$c_{1b}^{*}=-\log\lambda-\log(2\phi_{b})$。代入 $w=2\phi_{a}+\phi_{b}$ 得到
    \[
        \begin{aligned}
            c_{0}^{*}&=\frac{2\phi_{a}+\phi_{b}+\phi_{a}\log(2\phi_{a})+\phi_{b}\log(2\phi_{b})}{1+\phi_{a}+\phi_{b}}\\
            c_{1a}^{*}&=\frac{2\phi_{a}+\phi_{b}+\phi_{b}\log(2\phi_{b})-\log(2\phi_{a})-\phi_{b}\log(2\phi_{a})}{1+\phi_{a}+\phi_{b}}\\
            c_{1b}^{*}&=\frac{2\phi_{a}+\phi_{b}+\phi_{a}\log(2\phi_{a})-\log(2\phi_{b})-\phi_{a}\log(2\phi_{b})}{1+\phi_{a}+\phi_{b}}
        \end{aligned}
    \]
    \item 由一阶条件，$c_{1a}^{*}=c_{0}^{*}-\log(2\phi_{a})$，$c_{1b}^{*}=c_{0}^{*}-\log(2\phi_{b})$。记
    \[
        N=2\phi_{a}+\phi_{b}+\phi_{a}\log(2\phi_{a})+\phi_{b}\log(2\phi_{b}),\qquad
        D=1+\phi_{a}+\phi_{b},
    \]
    则 $c_{0}^{*}=N/D$。直接求导得
    \[
        \frac{\partial c_{0}^{*}}{\partial\phi_{a}}
        =\frac{3+\phi_{a}+\log(2\phi_{a})+\phi_{b}\bigl(2+\log(\phi_{a}/\phi_{b})\bigr)}{D^{2}},
        \qquad
        \frac{\partial c_{1a}^{*}}{\partial\phi_{a}}
        =\frac{\partial c_{0}^{*}}{\partial\phi_{a}}-\frac{1}{\phi_{a}},
    \]
    并且 $\partial c_{1b}^{*}/\partial\phi_{a}=\partial c_{0}^{*}/\partial\phi_{a}$。$\phi_{a}$ 上升时，禀赋价值 $w=2\phi_{a}+\phi_{b}$ 增加，同时状态 $a$ 的消费相对变贵；二者的相对强弱取决于价格，符号由分子决定。若 $2\phi_{a}\ge e^{-1}$ 且 $\phi_{a}\ge\phi_{b}$，分子为正，$c_{0}^{*}$ 与 $c_{1b}^{*}$ 上升；$c_{1a}^{*}$ 上升还要求 $\partial c_{0}^{*}/\partial\phi_{a}>1/\phi_{a}$。例如 $\phi_{a}=\phi_{b}=0.1$ 时分子为 $3.3+\log 0.2>0$，但导数约为 $1.17<10=1/\phi_{a}$，故 $c_{1a}^{*}$ 下降。若 $\phi_{a}=0.01$、$\phi_{b}=1$，分子为 $5.01+\log 0.02+\log 0.01<0$，$c_{0}^{*}$ 与 $c_{1b}^{*}$ 随 $\phi_{a}$ 上升而下降。

    对 $\phi_{b}$ 同样有
    \[
        \frac{\partial c_{0}^{*}}{\partial\phi_{b}}
        =\frac{2+\phi_{b}+\log(2\phi_{b})+\phi_{a}\log(\phi_{b}/\phi_{a})}{D^{2}},
        \qquad
        \frac{\partial c_{1b}^{*}}{\partial\phi_{b}}
        =\frac{\partial c_{0}^{*}}{\partial\phi_{b}}-\frac{1}{\phi_{b}},
    \]
    且 $\partial c_{1a}^{*}/\partial\phi_{b}=\partial c_{0}^{*}/\partial\phi_{b}$。若 $2\phi_{b}\ge e^{-1}$ 且 $\phi_{b}\ge\phi_{a}$，则 $c_{0}^{*}$ 与 $c_{1a}^{*}$ 不减。
    \item 当 $\phi_{a}=\phi_{b}=0.1$ 时，$w=0.3$，$D=1.2$，
    \[
        c_{0}^{*}=\frac{0.3+0.2\log(0.2)}{1.2}\approx -0.0182<0.
    \]
    该计划满足预算等式，但 $c_{0}^{*}<0$，因而不属于 (c) 的预算集。
\end{enumerate}
\end{solution}

\begin{problem}
\problemno{3.2} 考虑一个在 1 期只有一个可能状态的经济。（在这种情况下不存在不确定性。）参与者 1 的 0 期禀赋为 100 而 1 期禀赋为 1，即他的禀赋向量为 $\begin{bmatrix}100 \\ 1\end{bmatrix}$。他的偏好可以表示成如下形式：
\[
    U(c_{0},c_{1})=\log c_{0}+\rho\log c_{1}.
\]
有一只可交易证券，它的 0 期价格为 1、1 期支付为 $1+r_{F}$。这里，$r_{F}$ 是利率。
\begin{enumerate}[label=(\alph*)]
    \item 假设利率 $r_{F}$ 是给定的，导出参与者对证券的需求。
    \item 假设参与者 1 是经济中的唯一参与者。描述市场出清条件。
    \item 求解均衡利率。
    \item 均衡利率如何依赖于偏好参数 $\rho$？解释所得到的结论。
\end{enumerate}
\end{problem}
\begin{solution}
\begin{enumerate}[label=(\alph*)]
    \item 先不考虑消费的非负性，求解参与者的证券需求。此时参与者的优化问题变成了
    \[
        \max_{\theta}\bigl\{\log(100-\theta)+\rho\log\bigl(1+\theta(1+r_{F})\bigr)\bigr\}
    \]
    这与练习 2.7 是同一规划。一阶条件
    \[
        \frac{1}{100-\theta}=\rho\frac{1+r_{F}}{1+\theta(1+r_{F})}
    \]
    解得证券需求
    \[
        \theta^{*}=\frac{100\rho(1+r_{F})-1}{(1+\rho)(1+r_{F})}
    \]
    最优消费为
    \[
        c_{0}^{*}=\frac{100(1+r_{F})+1}{(1+\rho)(1+r_{F})},\quad
        c_{1}^{*}=\frac{\rho\bigl(100(1+r_{F})+1\bigr)}{1+\rho}
    \]
    由于最优消费都是非负的，因此上面的求解过程没有问题。
    \item 由于市场上只有一个参与者，因而市场出清条件是 $\theta^{*}=0$；
    \item 由上面的市场出清条件，我们可以求得均衡利率为 $r_{F}=\dfrac{1}{100\rho}-1$；
    \item $r_{F}$ 随着偏好系数的增加而减小。这是因为偏好系数增大，未来消费带来的边际效用相对增加，从而人们对能将财富在 0 期、1 期之间进行转换的证券需求增加，证券价格上升，利率下降。
\end{enumerate}
\end{solution}

\begin{problem}
\problemno{3.3*} 继续考虑练习 3.2 中定义的经济。现在假设还有另外一个参与者 2，他的效用函数和参与者 1 的一样但他的禀赋为 $\begin{bmatrix}1 \\ 100\end{bmatrix}$。
\begin{enumerate}[label=(\alph*)]
    \item 求解参与者 2 的最优消费/储蓄选择。
    \item 比较参与者 1 和 2 的最优消费/储蓄选择。
    \item 求解有两个参与者时的市场均衡。
    \item 计算均衡利率。它与前一情形（经济中只有参与者 1）下的均衡利率有何不同？
    \item 证明均衡配置是 Pareto 最优的。
\end{enumerate}
\end{problem}
\begin{solution}
\begin{enumerate}[label=(\alph*)]
    \item 先不考虑消费非负的约束，参与者 2 的优化问题就变成了
    \[
        \max_{\theta}\bigl\{\log(1-\theta)+\rho\log\bigl(100+\theta(1+r_{F})\bigr)\bigr\}
    \]
    一阶条件 $\dfrac{1}{1-\theta}=\rho\dfrac{1+r_{F}}{100+\theta(1+r_{F})}$ 给出证券需求
    \[
        \theta^{*}=\frac{\rho(1+r_{F})-100}{(1+\rho)(1+r_{F})}
    \]
    最优消费为
    \[
        c_{0}^{*}=\frac{(1+r_{F})+100}{(1+\rho)(1+r_{F})},\quad
        c_{1}^{*}=\frac{\rho\bigl((1+r_{F})+100\bigr)}{1+\rho}
    \]
    由于最优消费都是非负的，因此上面的求解过程没有问题。
    \item 两人的最优消费/储蓄选择形式上是完全一样的，差别在于两人的禀赋不一样；
    \item 此时的市场出清条件是 $\theta_{1}^{*}+\theta_{2}^{*}=0$，即
    \[
        \frac{100\rho(1+r_{F})-1}{(1+\rho)(1+r_{F})}+\frac{\rho(1+r_{F})-100}{(1+\rho)(1+r_{F})}=0
    \]
    分子之和为零：$101\rho(1+r_{F})=101$，故 $1+r_{F}=1/\rho$；
    \item 均衡利率是 $r_{F}=1/\rho-1$。只有参与者 1 时总禀赋为 $(100,1)$，均衡利率为 $1/(100\rho)-1$；现在总禀赋为 $(101,101)$，利率由跨期边际替代率 $\rho$ 单独决定，不再被极端的当期/未来禀赋扭曲；
    \item 配置 Pareto 最优，是指不存在另一种可行配置，使某人的效用严格提高而其他人都不下降。完全市场的均衡满足这一点（定理 3.4）；可检验的条件是两人的跨期边际替代率相等。边际效用之比为
    \[
        \frac{\partial U/\partial c_{1}}{\partial U/\partial c_{0}}=\rho\frac{c_{0}}{c_{1}}
    \]
    由两人的最优消费都有 $c_{0}^{*}/c_{1}^{*}=1/[\rho(1+r_{F})]$，因而两人的
    \[
        \rho\frac{c_{0}^{*}}{c_{1}^{*}}=\frac{1}{1+r_{F}}
    \]
    相等，均衡配置是 Pareto 最优的。
\end{enumerate}
\end{solution}

\begin{problem}
\problemno{3.4*} 在 1 期，经济有三个等可能的状态 $a$、$b$ 和 $c$。在证券市场上交易的 Arrow-Debreu 证券集合是完全的。经济中有两个参与者，参与者 1 和 2。他们的禀赋如下：
\begin{center}
\WJEventTree[left={$e_{1}{:}100$}]{0,0,0}
\qquad
\WJEventTree[left={$e_{2}{:}0$}]{200,100,50}
\end{center}
他们的偏好相同，形式如下：
\[
    U(c)=-\frac{1}{c_{0}}-\frac13\Bigl(\frac{1}{c_{1a}}+\frac{1}{c_{1b}}+\frac{1}{c_{1c}}\Bigr).
\]
\begin{enumerate}[label=(\alph*)]
    \item 求解每一参与者的最优消费/投资选择。
    \item 求解均衡状态价格。
    \item 证明对两个参与者的消费都有禀赋越高、消费越高的性质。
    \item 证明每一参与者的消费在总消费中所占的份额 $c_{k,1\omega}/e_{1\omega}$ 在所有状态下都是一样的，其中 $k=1,2$，而 $e_{1\omega}$ 是在状态 $\omega=a,b,c$ 下的 1 期总消费/禀赋。
\end{enumerate}
\end{problem}
\begin{solution}
\begin{enumerate}[label=(\alph*)]
    \item 求解最优消费选择时，先不考虑消费非负的约束。对于参与者 1，它的优化问题是：
    \[
        \begin{aligned}
            \max_{\{c_{0},c_{1}\}}\ &\Bigl\{-\frac{1}{c_{0}}-\frac13\Bigl(\frac{1}{c_{1a}}+\frac{1}{c_{1b}}+\frac{1}{c_{1c}}\Bigr)\Bigr\}\\
            \mathrm{s.t.}\quad &c_{0}+\phi_{a}c_{1a}+\phi_{b}c_{1b}+\phi_{c}c_{1c}=100
        \end{aligned}
    \]
    拉格朗日函数为
    \[
        \mathcal{L}=-\frac{1}{c_{0}}-\frac13\Bigl(\frac{1}{c_{1a}}+\frac{1}{c_{1b}}+\frac{1}{c_{1c}}\Bigr)+\lambda\bigl(100-c_{0}-\phi_{a}c_{1a}-\phi_{b}c_{1b}-\phi_{c}c_{1c}\bigr)
    \]
    一阶条件是边际效用与状态价格成比例：
    \[
        \frac{1}{c_{0}^{2}}=\lambda,\qquad
        \frac{1}{3c_{1\omega}^{2}}=\lambda\phi_{\omega},\quad\omega=a,b,c
    \]
    因而 $c_{1\omega}=c_{0}/\sqrt{3\phi_{\omega}}$。代入预算约束
    \[
        c_{0}\Bigl(1+\sqrt{\frac{\phi_{a}}{3}}+\sqrt{\frac{\phi_{b}}{3}}+\sqrt{\frac{\phi_{c}}{3}}\Bigr)=100
    \]
    得到参与者 1 的最优消费选择为
    \[
        \begin{aligned}
            c_{0}^{*}&=\frac{100}{1+\sqrt{\phi_{a}/3}+\sqrt{\phi_{b}/3}+\sqrt{\phi_{c}/3}}\\
            c_{1a}^{*}&=\frac{100}{\sqrt{3\phi_{a}}\bigl(1+\sqrt{\phi_{a}/3}+\sqrt{\phi_{b}/3}+\sqrt{\phi_{c}/3}\bigr)}\\
            c_{1b}^{*}&=\frac{100}{\sqrt{3\phi_{b}}\bigl(1+\sqrt{\phi_{a}/3}+\sqrt{\phi_{b}/3}+\sqrt{\phi_{c}/3}\bigr)}\\
            c_{1c}^{*}&=\frac{100}{\sqrt{3\phi_{c}}\bigl(1+\sqrt{\phi_{a}/3}+\sqrt{\phi_{b}/3}+\sqrt{\phi_{c}/3}\bigr)}
        \end{aligned}
    \]
    参与者 2 的财富为 $w_{2}=200\phi_{a}+100\phi_{b}+50\phi_{c}$，同一组一阶条件给出
    \[
        c_{0}^{*}=\frac{w_{2}}{1+\sqrt{\phi_{a}/3}+\sqrt{\phi_{b}/3}+\sqrt{\phi_{c}/3}},\qquad
        c_{1\omega}^{*}=\frac{c_{0}^{*}}{\sqrt{3\phi_{\omega}}}
    \]
    即
    \[
        \begin{aligned}
            c_{0}^{*}&=\frac{200\phi_{a}+100\phi_{b}+50\phi_{c}}{1+\sqrt{\phi_{a}/3}+\sqrt{\phi_{b}/3}+\sqrt{\phi_{c}/3}}\\
            c_{1a}^{*}&=\frac{200\phi_{a}+100\phi_{b}+50\phi_{c}}{\sqrt{3\phi_{a}}\bigl(1+\sqrt{\phi_{a}/3}+\sqrt{\phi_{b}/3}+\sqrt{\phi_{c}/3}\bigr)}\\
            c_{1b}^{*}&=\frac{200\phi_{a}+100\phi_{b}+50\phi_{c}}{\sqrt{3\phi_{b}}\bigl(1+\sqrt{\phi_{a}/3}+\sqrt{\phi_{b}/3}+\sqrt{\phi_{c}/3}\bigr)}\\
            c_{1c}^{*}&=\frac{200\phi_{a}+100\phi_{b}+50\phi_{c}}{\sqrt{3\phi_{c}}\bigl(1+\sqrt{\phi_{a}/3}+\sqrt{\phi_{b}/3}+\sqrt{\phi_{c}/3}\bigr)}
        \end{aligned}
    \]
    因为最优消费都是非负的，因而上面的求解过程是没有问题的；
    \item 市场出清条件是
    \[
        \begin{aligned}
            &\frac{100}{\sqrt{3\phi_{a}}\bigl(1+\sqrt{\phi_{a}/3}+\sqrt{\phi_{b}/3}+\sqrt{\phi_{c}/3}\bigr)}+\frac{200\phi_{a}+100\phi_{b}+50\phi_{c}}{\sqrt{3\phi_{a}}\bigl(1+\sqrt{\phi_{a}/3}+\sqrt{\phi_{b}/3}+\sqrt{\phi_{c}/3}\bigr)}=200\\
            &\frac{100}{\sqrt{3\phi_{b}}\bigl(1+\sqrt{\phi_{a}/3}+\sqrt{\phi_{b}/3}+\sqrt{\phi_{c}/3}\bigr)}+\frac{200\phi_{a}+100\phi_{b}+50\phi_{c}}{\sqrt{3\phi_{b}}\bigl(1+\sqrt{\phi_{a}/3}+\sqrt{\phi_{b}/3}+\sqrt{\phi_{c}/3}\bigr)}=100\\
            &\frac{100}{\sqrt{3\phi_{c}}\bigl(1+\sqrt{\phi_{a}/3}+\sqrt{\phi_{b}/3}+\sqrt{\phi_{c}/3}\bigr)}+\frac{200\phi_{a}+100\phi_{b}+50\phi_{c}}{\sqrt{3\phi_{c}}\bigl(1+\sqrt{\phi_{a}/3}+\sqrt{\phi_{b}/3}+\sqrt{\phi_{c}/3}\bigr)}=50
        \end{aligned}
    \]
    由 $c_{1\omega}=c_{0}/\sqrt{3\phi_{\omega}}$，0 期总禀赋为 100，故 $\sum_{k}c_{k,0}=100$。状态 $a$ 的出清化为
    \[
        \frac{\sum_{k}c_{k,0}}{\sqrt{3\phi_{a}}}=\frac{100}{\sqrt{3\phi_{a}}}=200
    \]
    即 $\sqrt{3\phi_{a}}=\dfrac12$，从而 $\phi_{a}=\dfrac{1}{12}$。同理 $\dfrac{100}{\sqrt{3\phi_{b}}}=100$ 给出 $\phi_{b}=\dfrac13$，$\dfrac{100}{\sqrt{3\phi_{c}}}=50$ 给出 $\phi_{c}=\dfrac43$。
    \item $\phi_{a}<\phi_{b}<\phi_{c}$，对于参与者 1、2 都有 $c_{1a}^{*}>c_{1b}^{*}>c_{1c}^{*}$，即总禀赋越高、消费越高；
    \item 因为对于消费者 1、2 都有
    \[
        \begin{aligned}
            \frac{c_{1a}^{*}}{c_{1b}^{*}}&=\sqrt{\frac{\phi_{b}}{\phi_{a}}}=2=\frac{200}{100}\\
            \frac{c_{1a}^{*}}{c_{1c}^{*}}&=\sqrt{\frac{\phi_{c}}{\phi_{a}}}=4=\frac{200}{50}
        \end{aligned}
    \]
    因此，每一参与者的消费在总消费中所占的份额在所有状态下都是相等的。
\end{enumerate}
\end{solution}

